% 3. Faza Algorytmiczna i Optymalizacja Ekstrakcji
% • Wybór najbardziej efektywnej metody post-processingu (tzw. wybielania) przy stałej
% liczbie źródeł (np. 64 RO).
% • Modyfikacja jednostki próbkowania (Sampling Unit) i implementacja wymiennych
% modułów:
% • XOR (rozwiązanie o najniższym koszcie sprzętowym).
% • CRC (standardowa metoda zaimplementowana w neoTRNG).
% • Toeplitz Extractor
% • AES-256 lub sha256
% • Pomiary wydajnościowe: przepustowość wyjściowa (Bitrate w Mbps), dodatkowe
% zużycie zasobów przez sam ekstraktor, pełna weryfikacja statystyczna NIST 800-22.
% • TABELA 2: Efektywność post-processingu i współczynnik $\eta$.
% • Zawartość: Typ ekstraktora | Dodatkowe LUT | Przepustowość [Mbps] |
% Wynik NIST 800-22 (Pass/Fail) | Współczynnik $\eta$ (Efektywność).

\chapter{Faza algorytmiczna i optymalizacja ekstrakcji}